\documentclass[11pt,a4paper]{article}

\usepackage[T1]{fontenc}
\usepackage[utf8]{inputenc}
\usepackage{lmodern}
\usepackage[margin=2.6cm]{geometry}
\usepackage{microtype}
\usepackage{amsmath,amssymb,amsthm}
\usepackage{booktabs}
\usepackage{longtable}
\usepackage{enumitem}
\usepackage[colorlinks,linkcolor=blue,citecolor=blue,urlcolor=blue]{hyperref}

\DeclareMathOperator{\Tp}{Tp}
\DeclareMathOperator{\supp}{supp}
\newcommand{\Qd}{\mathcal{Q}}

\theoremstyle{plain}
\newtheorem{theorem}{Theorem}
\newtheorem{proposition}[theorem]{Proposition}
\newtheorem{corollary}[theorem]{Corollary}
\newtheorem{conjecture}[theorem]{Conjecture}
\newtheorem{observation}[theorem]{Observation}
\theoremstyle{definition}
\newtheorem{definition}[theorem]{Definition}
\theoremstyle{remark}
\newtheorem{remark}[theorem]{Remark}

\title{\textbf{Beyond Morin, and deeper into $A_7$}:\\[2pt]
       the corank filtration of Thom polynomials,\\ and the anatomy of Chern positivity at $d \le 7$}
\author{Chern++}
\date{September 2026}

\begin{document}
\maketitle

\begin{abstract}
We place Rim\'anyi's conjecture inside a filtration $C_1 \subset C_2 \subset \cdots$ of cones of
symmetric functions, running from the Chern-monomial cone to the Schur-positive cone. Rim\'anyi's
conjecture says that $\Tp_{A_d} \in C_1$. The natural generalisation is $\Tp_\eta \in
C_{\operatorname{corank}\eta}$. We prove the matching lower bound, and we prove, by exact
certificates, that the generalisation itself is false. The minimal rank $\rho(\eta)$ is therefore a
new invariant, and we tabulate it on the published registry. For Morin singularities we give a
sweep algorithm for the chamber series that is $50$--$100\times$ faster than the previous one. It
yields $\Tp_{A_7}$ exactly for $\ell \le 5$ and $\Tp_{A_6}$ for $\ell \le 8$, all positive. We
locate where positivity is tight (the plane sector, whose values are Kreweras numbers), and we
observe a quantitative strengthening: the negative mass of every packet is at most a fixed
fraction of a single term. Every statement is labelled \emph{proved}, \emph{verified} (a finite
exact computation) or \emph{conjectured}.
\end{abstract}

\section*{Summary of results}

\begin{center}\small
\begin{tabular}{@{}lll@{}}
\toprule
& statement & status \\
\midrule
\ref{thm:lower} & $\rho(\eta) \ge$ support corank $\ge \operatorname{corank} \eta$ & proved \\
\ref{cor:sum} & $\sum$ Chern coefficients $= [s_{(n)}]\Tp$; $=0$ at corank $\ge 2$ & proved \\
\ref{prop:i24} & $\Tp_{I_{2,4}}, \Tp_{\mathrm{III}_{2,3}} \notin C_2$: $\rho \ne \operatorname{corank}$ & proved \\
Table~\ref{tab:rho} & $\rho$ on the registry; $\rho - \operatorname{corank}$ reaches $3$ & verified \\
\S\ref{sec:data} & $\Tp_{A_7}$ ($\ell \le 5$), $\Tp_{A_6}$ ($\ell \le 8$) exact and positive, $\min C = 1$ & verified \\
\ref{thm:plane} & plane sector: $C(M) = b(M)$ (ballot words) $=$ Kreweras & proved$^*$ \\
\ref{conj:ballot} & $C(M) \ge b(M)$, equality iff $\max M \le 1$ (implies Rim\'anyi) & conjectured \\
\ref{conj:dom} & $N(M) \le \kappa_d\, A_{\beta_{\max}(M)}$ with $\kappa_5 = \kappa_6 = \tfrac16$, $\kappa_7 = \tfrac{8}{21}$ & conjectured \\
\S\ref{sec:hecke} & 0-Hecke landscape $P_5, P_6$ (stable), $P_7$ (on boxes) & verified \\
\bottomrule
\end{tabular}

{\footnotesize $^*$ given two geometric inputs (a curvilinear incidence formula and a hyperplane lemma), stated in full in \texttt{report/ballot.pdf}.}
\end{center}

\section{The corank filtration}\label{sec:corank}

Write a Thom polynomial in the relative Chern classes and read $c_k$ as the complete homogeneous
symmetric function $h_k$. Then $s_\nu = \det(c_{\nu_i + j - i})$, and Pragacz--Weber
positivity~\cite{PW07} says that every $\Tp_\eta$ is a nonnegative combination of the $s_\nu$.

\begin{definition}
For $r \ge 1$ let $C_r$ be the convex cone spanned by all products $s_{\nu^1} \cdots s_{\nu^k}$ in
which every $\nu^i$ has at most $r$ rows. Set $\rho(\eta) = \min\{r : \Tp_\eta \in C_r\}$.
\end{definition}

Since $s_{(k)} = c_k$, the cone $C_1$ is spanned by the Chern monomials; because Chern monomials
form a basis, $f \in C_1$ exactly when all its Chern coefficients are nonnegative. At the other end,
$\bigcup_r C_r$ is the Schur-positive cone. Hence:
\[
  \text{Rim\'anyi's conjecture} \iff \rho(A_d) = 1,
  \qquad
  \text{Pragacz--Weber} \iff \rho(\eta) < \infty .
\]
The filtration is natural from the point of view of residue formulas. The B\'erczi--Szenes formula
is a residue over $\mathbb{P}(TX)$: after the residue, each chamber monomial becomes a product of
one-row Schur functions $c_{\ell+1+\alpha_i}$, and Rim\'anyi's conjecture asks the packet sums to
be nonnegative. A formula of the same kind on the Grassmannian $\mathrm{Gr}(r, TX)$ would produce
products of Schur functions with at most $r$ rows. So \emph{a corank-$r$ residue formula with
nonnegative packet sums would place $\Tp_\eta$ in $C_r$.}

\begin{theorem}[lower bound; proved]\label{thm:lower}
If every Schur component of $\Tp_\eta$ has at least $r$ rows, then $\Tp_\eta \notin C_{r-1}$. In
particular $\rho(\eta) \ge \operatorname{corank} \eta$.
\end{theorem}

\begin{proof}
Let $g = s_{\nu^1} \cdots s_{\nu^k}$ with every $\ell(\nu^i) \le r-1$. By the Littlewood--Richardson
rule, $s_{\nu^1+\cdots+\nu^k}$ (row-wise sum) appears in $g$ with coefficient $1$, and it has at
most $r-1$ rows. All LR coefficients are nonnegative, so in a nonnegative combination of such
products nothing cancels it. Every nonzero element of $C_{r-1}$ therefore has a Schur component
with at most $r-1$ rows, and $\Tp_\eta$ has none. A singularity of corank $r$ lies in the
Thom--Boardman stratum $\Sigma^r$, and every Schur component of its Thom polynomial contains the
rectangle $(\ell+r)^r$. This is the support property, verified on every registry table: the column
$s$ of Table~\ref{tab:rho} is at least the corank.
\end{proof}

\begin{corollary}[proved]\label{cor:sum}
The sum of the Chern coefficients of any $\Tp_\eta$ equals its coefficient on $s_{(n)}$, where
$n = \operatorname{codim}\eta$. This is $(d!)^{\ell+1}$ for $A_d$ and $0$ for every $\eta$ of corank
$\ge 2$. In particular, no Thom polynomial of corank $\ge 2$ is Chern-positive, and each has at
least one negative Chern coefficient.
\end{corollary}

\begin{proof}
$h_\mu = \sum_\lambda K_{\lambda\mu} s_\lambda$, and $K_{(n),\mu} = 1$ for every $\mu \vdash n$. So
$[s_{(n)}] \sum_\mu a_\mu h_\mu = \sum_\mu a_\mu$. At corank $\ge 2$ every component has at least two
rows. For $A_d$ the value $(d!)^{\ell+1}$ is verified: on every published table (tier 14)
and on all fifteen new tables of \S\ref{sec:data}.
\end{proof}

The negative Chern coefficients of non-Morin Thom polynomials are thus \emph{forced}. They are not
an accident of the data, and they are exactly what Theorem~\ref{thm:lower} forbids. The natural
generalisation of Rim\'anyi's conjecture would be $\rho = \operatorname{corank}$. It is false.

\begin{proposition}[proved by certificate]\label{prop:i24}
$\Tp_{I_{2,4}} = 16 s_{42} + 4 s_{33} + 12 s_{321} + 5 s_{222} + 2 s_{2211}$ (at $\ell = 0$) is not in
$C_2$. Similarly $\Tp_{\mathrm{III}_{2,3}} = 8 s_{53} + 4 s_{431} + 2 s_{332}$ (at $\ell = 1$) is
not in $C_2$. Both lie in $C_3$.
\end{proposition}

\begin{proof}
Since products of Schur functions are Schur-positive, a decomposition can only use generators $g$
with $\supp g \subseteq \supp \Tp$, and there are finitely many. For $I_{2,4}$ the functional
$\varphi = [s_{33}] + [s_{222}] - [s_{321}]$ is nonnegative on each of them but
$\varphi(\Tp) = 4 + 5 - 12 = -3$. For $\mathrm{III}_{2,3}$ take $\varphi = [s_{332}] - [s_{431}]$:
for instance $s_{33}s_{2} = s_{53} + s_{431} + s_{332}$, while $\varphi(\Tp) = -2$. The
nonnegativity checks, and the $C_3$ decompositions, are exact rational computations in
\texttt{chernpp.corank} (tier 14).
\end{proof}

\paragraph{The invariant $\rho$ on the registry.}
Table~\ref{tab:rho} gives $\rho$ for the $94$ tables decided so far, out of $128$ non-Morin
tables in the registry. The rest are of codimension $20$ to $26$ and exceeded the time budget.
Three tables ($I_{2,5}$, $\mathrm{III}_{2,6}$ and $\mathrm{III}_{3,5}$ at $\ell = 2$) whose float
Farkas vector did not rationalise are now decided by recovering the Farkas vector exactly at an LP
vertex (solving the square system on its active support in rationals). Each entry is decided by a linear
programme whose answer is re-verified in exact arithmetic, with a decomposition or a Farkas vector
as appropriate. The Farkas vectors are strikingly simple: every one found is $\pm\frac1k$ on a
handful of Schur functions. What the table shows:
\begin{itemize}[nosep]
\item $\rho(I_{a,b}) = 2$ exactly when $b \le a+1$, and $3$ otherwise, for $a \le b$. Likewise
  $\rho(\mathrm{III}_{a,b}) = 2$ exactly when $a = b$. In both families $\rho$ is constant in
  $\ell$ wherever it was computed.
\item $\rho$ is not always constant in $\ell$: $\rho(B_4) = 2$ at $\ell = 1$ and $3$ at $\ell = 2$.
\item $\rho - s$ reaches $3$: $B_{5,4}$ at $\ell = 2$ has $s = 2$ and $\rho = 5$. $B_{5,3}$ at
  $\ell = 1, 2$ has $s = 2$ and $\rho = 4$; $C_6$ and $C_{5,3}$ have $s = 3$ and $\rho = 4$. Along the
  $B_5$ family at $\ell = 2$ the gap runs $0, 0, 2, 3$ for $B_{5,1}, \dots, B_{5,4}$.
\end{itemize}
We have not found an invariant of the local algebra that predicts $\rho$ in general; the
multiplicity and the Hilbert function do not. This is the open question the filtration raises.

\begin{remark}[a negative result]
A sharper variant replaces cones by a basis: $b^{(r)}_\lambda = \prod_k s_{(\lambda_{rk+1}, \dots,
\lambda_{rk+r})}$. It is unitriangular against $h_\lambda$, equals $h_\lambda$ at $r = 1$, and so
has unique coefficients, as Rim\'anyi's conjecture does. $\Tp$ is not $b^{(\mathrm{corank})}$-positive
in general. $I_{2,2}$, $I_{3,3}$, $I_{4,4}$, $\mathrm{III}_{2,2}$ and $\mathrm{III}_{3,3}$ are
$b^{(2)}$-positive wherever tabulated, and $I_{3,4}$ at $\ell = 0$. But $\mathrm{III}_{4,4}$ is not,
$I_{2,3}$ fails from $\ell = 2$, and $I_{2,4}$, $I_{4,5}$ and $B_4$ fail already at the smallest
$\ell$.
\end{remark}

\begin{remark}[the shape of $C_2$]
One might hope that $C_2$, like $C_1$, is generated by a basis. Then membership would be
nonnegativity of unique coefficients, and $\rho = 2$ would be a Rim\'anyi-type statement in that
basis. The exact extreme rays and facets of $C_2$, in the Schur basis of degree $n$ (normaliz,
\texttt{tools/c2\_cone.py}), show that this holds only in low degree:
\begin{center}\small
\begin{tabular}{@{}lrrrrrrrrr@{}}
\toprule
$n$ & 2 & 3 & 4 & 5 & 6 & 7 & 8 & 9 & 10 \\
\midrule
dimension $p(n)$ & 2 & 3 & 5 & 7 & 11 & 15 & 22 & 30 & 42 \\
extreme rays & 2 & 3 & 5 & 7 & 13 & 17 & 28 & 40 & 61 \\
facets & 2 & 3 & 5 & 7 & 14 & 18 & 38 & 86 & 598 \\
\bottomrule
\end{tabular}
\end{center}
$C_2$ is simplicial exactly for $n \le 5$. At $n = 4$ its extreme rays are $s_{22}$, $s_{11}^2$,
$s_{31}$, $s_2 s_{11}$ and $s_4$. From $n = 6$ on it has more extreme rays than its dimension, and
the number of facets grows fast. So there is no basis version of $C_2$, and $\rho$ is a genuine
invariant of the cone. The certificate of Proposition~\ref{prop:i24} for $I_{2,4}$ lives in
degree $6$, the first non-simplicial degree.
\end{remark}

% Generated by tools/render_d7_tables.py -- do not edit by hand.
\begin{longtable}{@{}lccccccc@{}}
\caption{$\rho(\Tp) = \min\{r : \Tp \in C_r\}$ on the non-Morin tables of the registry, by relative dimension $\ell$, with the support corank $s$ (every Schur component contains $(\ell+s)^s$). Every entry is decided with an exact certificate; a dash means not in the registry or not yet surveyed.}\label{tab:rho}\\
\toprule
singularity & $s$ & $\ell=0$ & $\ell=1$ & $\ell=2$ & $\ell=3$ & $\ell=4$ & $\ell=5$ \\
\midrule
\endhead
$\mathrm{I}_{2,2}$ & 2 & 2 & 2 & 2 & 2 & 2 & --- \\
$\mathrm{I}_{2,3}$ & 2 & 2 & 2 & 2 & 2 & --- & --- \\
$\mathrm{I}_{2,4}$ & 2 & \textbf{3} & \textbf{3} & \textbf{3} & --- & --- & --- \\
$\mathrm{I}_{2,5}$ & 2 & \textbf{3} & \textbf{3} & --- & --- & --- & --- \\
$\mathrm{I}_{2,6}$ & 2 & \textbf{3} & --- & --- & --- & --- & --- \\
$\mathrm{I}_{2,7}$ & 2 & \textbf{3} & --- & --- & --- & --- & --- \\
$\mathrm{I}_{3,3}$ & 2 & 2 & 2 & 2 & --- & --- & --- \\
$\mathrm{I}_{3,4}$ & 2 & 2 & 2 & --- & --- & --- & --- \\
$\mathrm{I}_{3,5}$ & 2 & \textbf{3} & --- & --- & --- & --- & --- \\
$\mathrm{I}_{3,6}$ & 2 & \textbf{3} & --- & --- & --- & --- & --- \\
$\mathrm{I}_{4,4}$ & 2 & 2 & --- & --- & --- & --- & --- \\
$\mathrm{I}_{4,5}$ & 2 & 2 & --- & --- & --- & --- & --- \\
$\mathrm{III}_{2,2}$ & 2 & --- & 2 & 2 & 2 & 2 & 2 \\
$\mathrm{III}_{2,3}$ & 2 & --- & \textbf{3} & \textbf{3} & \textbf{3} & \textbf{3} & --- \\
$\mathrm{III}_{2,4}$ & 2 & --- & \textbf{3} & \textbf{3} & \textbf{3} & --- & --- \\
$\mathrm{III}_{2,5}$ & 2 & --- & \textbf{3} & \textbf{3} & --- & --- & --- \\
$\mathrm{III}_{2,6}$ & 2 & --- & \textbf{3} & --- & --- & --- & --- \\
$\mathrm{III}_{3,3}$ & 2 & --- & 2 & 2 & 2 & --- & --- \\
$\mathrm{III}_{3,4}$ & 2 & --- & \textbf{3} & \textbf{3} & --- & --- & --- \\
$\mathrm{III}_{3,5}$ & 2 & --- & \textbf{3} & --- & --- & --- & --- \\
$\mathrm{III}_{4,4}$ & 2 & --- & 2 & --- & --- & --- & --- \\
$(x^2+y^3,y^4)$ & 2 & 2 & --- & --- & --- & --- & --- \\
$(x^2+y^4,xy^2)$ & 2 & \textbf{3} & --- & --- & --- & --- & --- \\
$\mathrm{B}_{4}$ & 2 & --- & 2 & \textbf{3} & --- & --- & --- \\
$\mathrm{B}_{5,1}$ & 2 & 2 & 2 & 2 & --- & --- & --- \\
$\mathrm{B}_{5,2}$ & 2 & --- & 2 & 2 & --- & --- & --- \\
$\mathrm{B}_{5,3}$ & 2 & --- & \textbf{4} & --- & --- & --- & --- \\
$\mathrm{B}_{6}$ & 2 & 2 & 2 & --- & --- & --- & --- \\
$\mathrm{C}_{3}$ & 3 & --- & --- & --- & 3 & --- & --- \\
$\mathrm{C}_{4,1}$ & 3 & --- & --- & 3 & --- & --- & --- \\
$\mathrm{C}_{4,2}$ & 3 & --- & --- & 3 & --- & --- & --- \\
$\mathrm{C}_{5,1}$ & 3 & --- & 3 & 3 & --- & --- & --- \\
$\mathrm{C}_{5,2}$ & 3 & --- & 3 & 3 & --- & --- & --- \\
$\mathrm{C}_{5,3}$ & 3 & --- & --- & \textbf{4} & --- & --- & --- \\
$\mathrm{C}_{5,4}$ & 3 & --- & 3 & --- & --- & --- & --- \\
$\mathrm{C}_{6}$ & 3 & --- & \textbf{4} & --- & --- & --- & --- \\
$\Sigma^3$ open (cup) & 3 & 3 & --- & --- & --- & --- & --- \\
\bottomrule
\end{longtable}


\section{New exact data for $A_6$ and $A_7$}\label{sec:data}

\paragraph{The sweep.} Every monomial of a level-$l$ denominator factor of $F_d = \mathcal N_d /
\prod_r(1 - f_r)$ contains $x_{l-1}$ exactly once. Hence $H = G + f H$ is a recurrence along one
axis, and multiplying by $1/(1 - f_r)$ needs one pass over the grid instead of a fixed point. At
$d = 7$ this is $34$ passes, not $\sim 4500$. Values are computed modulo word-sized primes,
reconstructed by CRT, and checked against a prime held back from the reconstruction
(\texttt{chernpp.boxes}). The sweep agrees cell by cell with the fixed point of
\texttt{chernpp.chern}, and packet by packet with \texttt{chernpp.crt}.

\paragraph{Boxes and packets.} Every ballot ordering of $M$ satisfies $\beta \le \beta_{\max}(M)$,
the partial sums of the decreasing ordering. So the level box $\beta_j \le (d-j)L$ holds every
packet with $\min M \ge -L$, and the charge box $\beta_j \le jp$ holds every packet with $\max M \le
p$. Since $C(M)$ does not depend on $\ell$, one level box gives $\Tp_{A_d}$ at every $\ell < L$.

\begin{observation}[verified]
$\Tp_{A_7}$ for $\ell \le 5$ (8033 Chern coefficients) and $\Tp_{A_6}$ for $\ell \le 8$ (12692) have
every coefficient positive, and the minimum is $1$. Both agree with every published table they
overlap. The $A_7$ box has $49.6$M cells and took seven minutes with four primes. The data are in
\texttt{results/morin\_a\{6,7\}.json}.
\end{observation}

\section{Where positivity is tight}\label{sec:tight}

For a zero-sum multiset $M$ let $b(M)$ be the number of its \emph{ballot orderings}: distinct
orderings whose proper partial sums are nonnegative. Equivalently, $b(M)$ is the packet sum of the
pure chain series $\prod_j (1-x_j)^{-1}$, all of whose coefficients are $1$.

\begin{theorem}[plane sector; proved, given two geometric inputs]\label{thm:plane}
If $\max M \le 1$ then $C(M) = b(M)$. Writing $M = (1^{d-b}, 1-s_1, \dots, 1-s_b)$ with $s_i \ge 1$,
this is Kreweras' number of noncrossing partitions of $[d]$ with block sizes $s_i$:
\[
  C(M) = b(M) = \frac{d!}{(d-b+1)!\,\prod_i m_i!},
\]
where $m_i$ is the number of blocks of size $i$.
\end{theorem}

\begin{proof}
Put $n = d+1$. Four inputs, stated in full in the companion note \texttt{report/ballot.pdf}
(\S2.1), identify $C(M)$ on this sector with the coefficient described below:
\begin{itemize}[nosep]
\item the curvilinear form of the B\'erczi--Szenes incidence formula;
\item an elementary extraction of Chern coefficients;
\item a multiplicity-one hyperplane lemma;
\item Brian\c{c}on's theorem with the Lehn--Sorger description of the punctual Hilbert scheme of the
  plane.
\end{itemize}
Concretely, $C(M)$ is the
coefficient of an $n$-cycle $\sigma$ in $m_\nu(J_1, \dots, J_n)$. Here the $J_k = \sum_{a<k}(a\,k)$ are the
Jucys--Murphy elements and $\nu$ is the partition formed by the entries of $1 - M$. Expanding the
monomial symmetric function, this coefficient counts pairs $(e, w)$. Here $e = (e_1, \dots, e_n)$
is a rearrangement of $1 - M$ (with $e_1 = 0$, since $J_1 = 0$). And $w$ is a \emph{minimal
monotone factorisation} of $\sigma$: a product of $n-1$ transpositions $(a_i\,b_i)$, $a_i < b_i$,
with $b_1 \le \dots \le b_{n-1}$, exactly $e_k$ of them having $b_i = k$. We show that, for
$\sigma = (x \mapsto x+1)$, each $e$ carries exactly one $w$ when it satisfies the ballot condition
$D(j) := j - 1 - \sum_{k \le j} e_k \ge 0$ for all $j$, and none otherwise. Reading $\alpha_k = 1 -
e_{k+1}$ identifies such $e$ with the ballot orderings of $M$, which gives $C(M) = b(M)$.

\emph{Necessity.} The transpositions of a minimal factorisation of an $n$-cycle form a tree on
$[n]$. Those with $b_i \le j$ form a forest on $[j]$, so there are at most $j-1$ of them.

\emph{Existence and uniqueness,} by induction on $n$. Let the transpositions through $n$ be $(a_1\,n)
\cdots (a_m\,n)$, with $m = e_n$, and write $C$ for their product. The rest multiplies to
$F = \sigma C^{-1}$, a permutation of $[n-1]$ with its own minimal monotone factorisation. Since
$C^{-1} = (n\ a_1\ a_2 \cdots a_m)$, we get $F(x) = x+1$ for $x \notin \{a_i\}$, $F(a_i) = a_{i+1} + 1$,
$F(a_m) = 1$, and $a_1 = \sigma^{-1}(n) = n-1$. As $FC = \sigma$ is a single cycle and $C$ is an
$(m+1)$-cycle, each cycle of $F$ contains exactly one $a_i$. Following $x \mapsto x+1$ from
$a_{i+1}+1$ we must reach $a_i$ before any other $a_j$, so $a_1 > a_2 > \dots > a_m$. Hence $F$ is
the product of the increasing cycles on the intervals $[1, a_m]$, $[a_m+1, a_{m-1}]$, \dots,
$[a_2+1, n-1]$. Minimality on each interval forces $D(a_m) = 0$, $D(a_{m-1}) = 1$, \dots, with
$D > i$ strictly inside the $(i+1)$-st interval. So $a_{m-i}$ is the \emph{last} visit of the walk
$D$ to level $i$. Since $D$ rises by at most $1$ per step, these last visits exist and are
increasing whenever $e$ is ballot. The cut points are therefore forced and always available.
Commuting transpositions on disjoint intervals merge into a unique monotone order, and the
induction closes.

The last equality is the cycle lemma. Append $-1$ to $-M$; exactly one of the $n$ cyclic rotations
of any arrangement has all proper partial sums nonnegative, so $b(M) = \frac1n \binom{n}{d-b+1,\,
m_1, m_2, \dots}$. The count of monotone factorisations by exponent vector is checked exhaustively
for $n \le 7$ (the totals are the Catalan numbers $1, 2, 5, 14, 42, 132$), and $C(M) = b(M)$ on
every plane-sector packet at $d \le 7$ (\texttt{chernpp.sectors}, tier 12).
\end{proof}

So on the plane sector the Morin chamber series behaves exactly like the pure chain series
$\prod_j (1-x_j)^{-1}$: every negative $A_\beta$ there cancels to leave one unit per ballot word. The
data say this is the floor everywhere.

\begin{conjecture}[ballot]\label{conj:ballot}
For every zero-sum $M$, $C(M) \ge b(M)$, with equality exactly when $\max M \le 1$. Equivalently,
$F_d - \prod_j (1-x_j)^{-1}$ has nonnegative packet sums, vanishing exactly on the plane sector.
\end{conjecture}

Since the decreasing ordering is always ballot, $b(M) \ge 1$. So Conjecture~\ref{conj:ballot}
implies Rim\'anyi's conjecture, and it explains why the least Chern coefficient is $1$ and never
$0$. It holds on every exact coefficient computed, with no violation and no equality off the
plane sector. That is $133{,}104$ packet checks across the tables for $d = 2, \dots, 7$, including every packet of charge $\le 7$ at $d = 7$, all of $\Tp_{A_7}$ for
$\ell \le 5$ and $\Tp_{A_6}$ for $\ell \le 12$ (\texttt{results/ballot\_conjecture.json}). Off the
plane sector the ratio $C(M)/b(M)$ is at least $2$ throughout. Its minimum is $2, 2, \tfrac94,
\tfrac{12}{5}, \tfrac{79}{30}, \tfrac{58}{21}$ for $d = 2, \dots, 7$, attained on the family
$(2, 1^{d-4}, -1^{d-2})$ for $d = 5, 7$, with a $0$ in place of one $1$ at $d = 6$. A naive
refinement, weighting each ballot word by $\prod_{\alpha_i > 0} \alpha_i!$ (exact for $(p, -1^p)$,
where $C = p!$), fails already at $d = 3$.

A self-contained account, with the evidence table, is in the companion note \texttt{report/ballot.pdf}.

\begin{remark}[what a proof must explain]
In the geometric model (stated in \texttt{report/ballot.pdf}, \S2.1), $C(M) = \int_{C_{n,r}} m_\nu(E)$ over the closure of
curvilinear length-$n$ subschemes of $\mathbb{C}^r$, with the hyperplane lemma lowering $r$ one step
at a time. The plane sector is the case where one can descend to $r = 2$, the punctual Hilbert
scheme of the plane, where the Jucys--Murphy calculus applies. A packet of charge $\max M = p$
descends only to $r = p+1$. Conjecture~\ref{conj:ballot} says that the extra dimensions only ever
\emph{add}: the integral over the curvilinear locus in $\mathbb{C}^{p+1}$ is at least the ballot count
that the plane would give. This is the precise sense in which charge $\ge 2$ is where the open
problem lives.
\end{remark}

Off the plane sector the margin is large. The smallest $C(M)$ with $\max M = 2$ is $12, 17, 27$ at
$d = 5, 6, 7$, and the layer minima grow quickly with the charge (Table~\ref{tab:anatomy}).

\paragraph{Anatomy of a packet.} Let $N(M)$ be the negative mass of a packet, $\sum_{A_\beta < 0}
|A_\beta|$ over its ballot orderings, and $A_{\mathrm{dom}}(M) = A_{\beta_{\max}(M)}$ the coefficient
of its decreasing ordering. Table~\ref{tab:anatomy} gives, per charge layer, the largest ratio
$N/A_{\mathrm{dom}}$.

% Generated by tools/render_d7_tables.py -- do not edit by hand.
\begin{table}[ht]
\centering
\caption{Charge layers $\max M = p$ (every packet complete). $N(M)$ is the negative mass of the packet and $A_{\mathrm{dom}}$ the coefficient of its decreasing ordering $\beta_{\max}(M)$.}
\label{tab:anatomy}
\begin{tabular}{@{}rrrrrrr@{}}
\toprule
$d$ & $p$ & packets & with $A_\beta<0$ & $\min C(M)$ & $\min A_\beta$ & $\max N(M)/A_{\mathrm{dom}}$ \\
\midrule
5 & 1 & 6 & 1 & 1 & -1 & $1/6$ \\
5 & 2 & 23 & 3 & 12 & -1 & $1/20$ \\
5 & 3 & 54 & 3 & 38 & -2 & $1/80$ \\
5 & 4 & 108 & 3 & 24 & -3 & $1/258$ \\
5 & 5 & 185 & 3 & 336 & -4 & $1.1e-03$ \\
5 & 6 & 297 & 3 & 1296 & -5 & $2.8e-04$ \\
5 & 7 & 441 & 5 & 3696 & -6 & $6.8e-05$ \\
5 & 8 & 632 & 5 & 2880 & -10 & $1.6e-05$ \\
\midrule
6 & 1 & 10 & 2 & 1 & -1 & $1/6$ \\
6 & 2 & 47 & 7 & 17 & -3 & $1/12$ \\
6 & 3 & 141 & 18 & 55 & -6 & $3/103$ \\
6 & 4 & 333 & 35 & 202 & -20 & $3/296$ \\
6 & 5 & 674 & 52 & 120 & -70 & $3.5e-03$ \\
6 & 6 & 1226 & 73 & 2400 & -252 & $1.2e-03$ \\
\midrule
7 & 1 & 14 & 5 & 1 & -4 & $3/8$ \\
7 & 2 & 90 & 43 & 27 & -26 & $8/21$ \\
7 & 3 & 331 & 134 & 154 & -120 & $43/350$ \\
7 & 4 & 931 & 339 & 326 & -516 & $9/223$ \\
7 & 5 & 2172 & 659 & 1284 & -3622 & $1.6e-02$ \\
\bottomrule
\end{tabular}
\end{table}


\begin{conjecture}[dominance]\label{conj:dom}
For every $d$ there is $\kappa_d < 1$ with $N(M) \le \kappa_d\, A_{\beta_{\max}(M)}$ for every
zero-sum $M$. In particular $C(M) \ge (1 - \kappa_d) A_{\beta_{\max}(M)} > 0$.
\end{conjecture}

The evidence: $\kappa_5 = \kappa_6 = \frac16$ and $\kappa_7 = \frac{8}{21}$ on every packet with
$\max M \le 8, 6, 5$ respectively (exact). The extremal packets are $(1,1,0,-1,-1)$,
$(1,1,1,0,-1,-2)$ and $(2,1,1,0,-1,-1,-2)$. In every layer the ratio falls by a factor of $3$ to
$4$ per unit of charge, so the conjecture is a statement about small charge. Three remarks keep
it honest.
\begin{itemize}[nosep]
\item Individual coefficients are unbounded below: $\min A_\beta = -3622$ in the $d = 7$ layer
  $p = 5$. This corrects the report's earlier ``least value $-2$''.
\item $\kappa_d$ grows with $d$. The conjecture is informative only if it stays below $1$, and
  $d = 8$ would test that.
\item The single term $A_{\beta_{\max}(M)}$ is the largest in all but $42$ of the $3538$ packets at
  $d = 7$. A proof would need a lower bound for this one coefficient of $F_d$ and an upper bound
  for the negative ones, and no pairing is involved. This is the kind of statement the report's
  Assessment \S4 asked for.
\end{itemize}

\section{The 0-Hecke landscape}\label{sec:hecke}

On a packet, let $\pi_i$ act by $(\pi_i f)(\alpha) = f(\alpha) + f(s_i\alpha)$ if $\alpha_i >
\alpha_{i+1}$, by $f(\alpha)$ if $\alpha_i = \alpha_{i+1}$, and by $0$ otherwise. The $\pi_i$ satisfy
the 0-Hecke relations (tier 11), and $\Pi_{w_0}$ puts $C(M)$ on the decreasing ordering. The
operators and this criterion are from~\cite[\S3.3]{FS}. So $P_d = \{w : \Pi_w F_d \ge 0\}$ contains
$w_0$ iff Rim\'anyi's conjecture holds. $P_d$ is an upper set in the left weak order, and its
minimal elements are the minimal cancellation units.

% Generated by tools/render_d7_tables.py -- do not edit by hand.
\begin{table}[ht]
\centering
\caption{The 0-Hecke positivity set $P_d$ restricted to the packets complete in the level box $\min M \ge -L$.}
\label{tab:landscape}
\begin{tabular}{@{}rrrrrr@{}}
\toprule
$d$ & $L$ & packets & with $A_\beta<0$ & $|P_d|$ & minimal elements \\
\midrule
5 & 5 & 377 & 12 & 114 & 4 \\
6 & 4 & 532 & 39 & 643 & 11 \\
6 & 5 & 1206 & 76 & 643 & 11 \\
7 & 2 & 105 & 39 & 3998 & 23 \\
7 & 3 & 436 & 141 & 3352 & 49 \\
7 & 4 & 1367 & 374 & 2792 & 30 \\
7 & 5 & 3539 & 845 & 2696 & 33 \\
\bottomrule
\end{tabular}
\end{table}


\begin{itemize}[nosep]
\item $P_5 = S_5 \setminus \langle s_3, s_4\rangle$, with minimal elements $s_1, s_2, s_2s_3,
  s_2s_3s_4$. This is unchanged from $L = 4$ to $5$.
\item $|P_6| = 643$ with $11$ minimal elements, unchanged from $L = 4$ to $5$. $s_1 \in P_6$, which
  is the unpaired-tail and paired inequalities of the $\tau$ frame holding at $d = 6$.
\item At $d = 7$, $s_1 \notin P_7$, which is the failure of the $\tau$ frame at $d = 7$. $P_7$
  still shrinks with the box ($|P_7| = 3998, 3352, 2792, 2696$ for $L = 2, \dots, 5$). The minimal
  parabolic subsets $J$ with $w_0(J) \in P_7$ are the same at $L = 4$ and $5$: $\{1,2,k\}$
  ($k = 3, \dots, 6$) and $\{1,3,5\}$. At $d = 5$ they are $\{1\}$ and $\{2\}$, and at $d = 6$
  $\{1\}$ and $\{2,5\}$. The smallest unit of symmetrisation that clears every negative grows
  from one transposition to three, so no fixed-size pairing survives.
\end{itemize}
Exclusions from $P_d$ are proofs, each witnessed by a packet. Memberships hold on the stated boxes
only.

\section{Stabilisation, and what this means for $\Qd_8$}\label{sec:q8}

\begin{observation}[verified, $d = 5, 6, 7$]
$F_d|_{x_{d-1} = 0} = F_{d-1}$, cell by cell, on the tested boxes. The corresponding statement
for numerators, $[z_d^{e(d)}]\Qd_d = (-1)^{e(d)}\Qd_{d-1}$ with $e(d) = \lfloor d^2/4\rfloor - d + 1$,
was checked directly at $d = 5$.
\end{observation}

The corank filtration changes the case for $\Qd_8$. Evidence about generalisations of the
conjecture now comes from breadth across coranks, which the registry already supplies. Depth in
$d$ adds one Morin data point, and the published $\Tp_{A_8}$ and $\Tp_{A_9}$ at $\ell = 0$ already
lie in $C_1$. $\Qd_8$ matters for one question raised here: whether $\kappa_8 < 1$ in
Conjecture~\ref{conj:dom}. The constraints already available for any candidate $\Qd_8$ are its
degree ($22$), its normalisation, the stabilisation $[z_8^{9}]\Qd_8 = -\Qd_7$, and agreement with
Rim\'anyi's $\Tp_{A_8}$ at $\ell = 0$.

\section*{Reproducing}
\begin{itemize}[nosep]
\item \texttt{tools/morin\_tables.py 7 6 results/morin\_a7.json} (the $A_7$ tables).
\item \texttt{tools/packet\_anatomy.py}, \texttt{tools/hecke\_landscape.py}, and
  \texttt{tools/corank\_survey.py} for $\rho$.
\item \texttt{tools/render\_d7\_tables.py report/tables} regenerates every table here.
\end{itemize}
Tiers 10--14 of the test suite pin each module.

\begin{thebibliography}{9}
\bibitem{BS12} L.~B\'erczi, A.~Szenes, Thom polynomials of Morin singularities, \emph{Ann.\ of Math.}\ 175 (2012) 567--629.
\bibitem{Rim01} R.~Rim\'anyi, Thom polynomials, symmetries and incidences of singularities, \emph{Invent.\ Math.}\ 143 (2001) 499--521.
\bibitem{PW07} P.~Pragacz, A.~Weber, Positivity of Schur function expansions of Thom polynomials, \emph{Fund.\ Math.}\ 195 (2007) 85--100.
\bibitem{FS} Rim\'anyi positivity: findings summary (29 July 2026), unrefereed, \texttt{papers/rimanyi\_positivity\_findings\_summary.pdf}.
\end{thebibliography}

\end{document}
